\section{From Embryonic Neurogenesis to Postnatal and Adult Neurogenic Processes}
Toward the end of neurogenesis ($\sim$E16.5 onward), progenitor cells transition from neurogenesis to gliogenesis, giving rise to distinct glial precursors that subsequently differentiate into astrocytes and oligodendrocytes \citep{VIEIRANeuralStemCell2018}.
In the present work, however, we focus exclusively on neurogenic processes and do not further address gliogenesis.

\section{Species-Specific Differences: From Rodent Models to Human Brain Development}
With respect to rodents such as mice and rats, which are widely used as laboratory models, it is important to note that most mammals possess larger and more complex brains. This increased complexity is reflected in greater cellular diversity, additional cortical layers, and fundamental differences in developmental dynamics. For instance, neurogenesis in the human cortex is markedly prolonged, lasting approximately 141 days compared to around 11 days in mice, and is accompanied by an extended proliferative phase and longer cell cycle durations in human NPCs \citep{LIBE-PHILIPPOTCellularMolecularMechanisms2021}.
\\

Structural differences further highlight this divergence. The proliferative SVZ is substantially expanded in humans and subdivided into an inner and outer SVZ. Within the outer SVZ, BPs are greatly enriched and have been closely associated with increased brain size and cortical expansion \citep{reilloRoleIntermediateRadial2011}.
In addition, BPs in humans exhibit enhanced proliferative capacity compared to their rodent counterparts. Finally, these developmental and cellular differences are paralleled by large-scale anatomical distinctions, as mice are lissencephalic, possessing smooth brains, whereas humans exhibit gyrencephaly, characterized by varying degrees of cortical folding \citep{BORRELLHowCellsFold2018}.
\\

Nevertheless, the cerebral cortex of rodents and humans share fundamental structural and functional properties, including layered organization, and similar developmental sequences \citep{luhmannCanWeUnderstand2020}, that allows detailed \textit{in vivo} and \textit{in vitro} analyses, as well as specific experimental manipulations.

\section{Developmental Dynamics during Neurogenesis}
Development is not a static process, but a finely tuned interplay of coordinated dynamics, that guide cell behavior at every step. By dynamics, we mean the continuous changes in a cell’s internal state, including its transcriptome, proteome, and metabolome, as well as the evolving external milieu of signaling molecules, mechanical forces and environmental cues.
\\

Together, these dynamics shape the structure of the cortex. They drive the transition of NSCs and NPCs from multipotency to lineage-restricted states, ultimately giving rise to specialized neural and glial identities. Additionally, they regulate the balance between self-renewal and differentiative state.

\subsection{Temporal Remodeling of the Transcriptome}
Advances in high-resolution molecular profiling have begun to transform our understanding of these developmental processes by enabling the characterization of cellular states at a remarkably detailed view (reviewed in \citep{BRISCOELookingNeurodevelopmentBig2020}).  In particular, single-cell RNA sequencing (\gls{scRNA-seq})
and spatial transcriptomic approaches have substantially refined our understanding of individual cell types and developmental trajectories throughout cortical development. For example, these technologies have revealed that NPCs and their progeny transition through a series of transient molecular states, characterized by sequential activation and repression of gene networks. This dynamic transition of transcriptional profiles not only distinguishes emerging cell types, but also occurs within individual progenitor populations \citep{RUANProgenitorCellDiversity2021}, where cells exhibit asynchronous progression and diverse lineage trajectories. 
\\

The identity of each cell type is defined by characteristic marker gene expression patterns, yet during neurogenesis the transcriptome undergoes extensive and continuous remodeling, which complicates the characterization of cell types during embryonic development. Early embryonic mouse brain studies have demonstrated that as much as 10–13\% of mRNAs and proteins
change in expression within a two-day developmental window \citep{HARTLTranscriptomeProteomeAnalysis2008}. At E9.5, gene expression is dominated by transcripts associated with metabolism and the cell cycle, reflecting active proliferation. By E11.5 and E13.5, coinciding with the transition from progenitor states to neuronal differentiation, there is a marked shift toward the expression of genes involved in neuronal identity and function.
\\

Consistent with this view, scRNA-seq analyses of human fetal brain tissue have revealed that neural progenitor populations are not only heterogeneous but also dynamically diverge along distinct developmental trajectories \citep{SILVESTRISSinglecellTranscriptomicProfiling}. For example, BPs from the SVZ segregate into subpopulations that follow two principal transcriptional programs, leading either toward glial lineages or toward neuronal identities. Similar observations from in vitro systems further support the idea that multiple progenitor subpopulations can coexist while undergoing progressive transcriptional shifts.
\\

Importantly, transcriptional dynamics are not restricted to embryonic development but persist across the lifespan. Spatial and temporal transcriptomic studies in the mouse brain have shown that gene expression profiles continue to evolve from early postnatal stages into adulthood and aging. Early postnatal brains are enriched for transcriptional programs associated with neurogenesis, gliogenesis, and myelination, whereas aged brains exhibit increased expression of genes linked to stress responses and inflammation \citep{CONACHERSpatialTranscriptomicsReveals2025}.

\section{Overview of Cellular Energy Metabolism}
\begin{figure}[H]
    \centering
    \includegraphics[width=\textwidth]{mitoMammal_model_scheme_Maelle.png}
    \caption{\textbf{ Schematic overview of cellular energy metabolism.} Overview of key reactions within the MitoMammal \citep{CHAPMANMitoMAMMALGenomeScale2025}  model, consisting of 784 genes (human and mouse orthologs) encoding 560 reactions. Figure created with BioRender.com.
    }
    \label{fig:mitomammal}
\end{figure}

Energy requirements in the brain change substantially over the course of development. In adulthood, the human brain consumes approximately 15-20\%  of the body’s resting metabolic resources, despite accounting for only about 2\% of total biomass \citep{HOFMANEnergyMetabolismBrain1983}.
In contrast, during early childhood, glucose utilization is markedly elevated, with uptake per gram of tissue peaking at roughly twice of the adult levels around the age of two to three years \citep{CHUGANIPositronEmissionTomography1987}.
\\

Our understanding of energy metabolism during embryonic brain development remains limited. Direct \textit{in vivo}  measurements are currently not feasible due to both technical and ethical constraints. Established approaches used to quantify postnatal glucose metabolism such as 2-deoxyglucose uptake assays or Fludeoxyglucose-Positron emission tomography (\gls{FDG-PET}) imaging, cannot
be readily applied \textit{in utero}, given the small size and inaccessibility of the embryonic neocortex for live imaging. As a result, the precise energetic demands of the developing neocortex are still not well defined.
\\

At the cellular level, these energetic demands are primarily met through the breakdown of carbon substrates, predominantly glucose in the brain, to generate adenosine-5’-triphosphate( \gls{ATP}),  which serves as the universal energy currency of the cell. ATP is generated predominantly through two major metabolic pathways: glycolysis in the cytosol and oxidative phosphorylation (\gls{OXPHOS}) in the mitochondria.

\subsection{Glycolysis}
Glycolysis derives its name from the Greek \textit{glykys} (“sweet”) and \textit{lysis} 
(“splitting”), reflecting the breakdown of one glucose molecule into two molecules of pyruvate in the cytosol. It consists of a ten-step enzymatic pathway driven by key enzymes such as hexokinase (\gls{HK}), 
phosphofructokinase-1 (\gls{PFK1}), phosphoglycerate kinase (\gls{PGK}) and pyruvate kinase (\gls{PYK}). The latter two enzymes are responsible for the anaerobic ATP production (see Fig~\ref{fig:mitomammal} and Fig~\ref{fig:glycolysis_TCA} A)
\\

Glucose is first converted to glucose-6-phosphate by HK, and subsequently isomerized to fructose-6-phosphate by phosphoglucose isomerase (\gls{GPI}).
By using one ATP molecule, PFK1 produces fructose-1,6-bisphosphate. Fructose bisphosphate aldolase then catalyzes the split of fructose-1,6-bisphosphate into dihydroxyacetone phosphate (\gls{DHAP})
and glyceraldehyde 3-phosphate (\gls{G3P}), where DHAP will be converted to G3P by triosephosphate isomerase (\gls{TPI1}).
Now, 2 molecules of G3P continue glycolysis with the oxidation into 1,3-bisphosphoglycerate while reducing one NAD$^+$ molecule to NADH each. Next, PGK catalyzes a phosphorylation step, in which 13DPG is converted to 3-phosphoglycerate , alongside with the production of one ATP molecule. 3-phosphoglycerate is then converted into 2-phosphoglycerate via phosphoglycerate mutase (\gls{PGAM}),
, which is further converted into phosphoenolpyruvate (\gls{PEP}) by enolase. In the final step of glycolysis, PEP donates its phosphate group to generate ATP, forming pyruvate through the action of PYK. 
\\

Importantly, glycolysis does not require oxygen and can operate under both aerobic and anaerobic conditions. Even when oxygen is available, cells may continue to rely on glycolysis without subsequent mitochondrial oxidation, a phenomenon known from cancer cells and often referred to as aerobic glycolysis or the Warburg effect \citep{THOMPSONCenturyWarburgEffect2023}.
\\

Pyruvate, the end-product of glycolysis, serves as a central metabolic branching point, as it can be directed into multiple downstream pathways. For instance, it may be reduced to lactate by lactate dehydrogenase A \gls{LDHA}
during fermentation. Alternatively, under aerobic conditions, pyruvate is typically transported into the mitochondria, where it is oxidized into acetyl-coenzyme A (\gls{CoA}),  which then enters the tricarboxylic acid (\gls{TCA}) cycle (see Fig~\ref{fig:glycolysis_TCA} B).

\begin{figure}[H]
    \centering
    \includegraphics[width=\textwidth]{Glycolysis_TCA.png}
    \caption{\textbf{Glycolysis and TCA cycle.}
    Overview of the reactions involved in \textbf{(A)} glycolysis and \textbf{(B)} tricarboxylic acid (TCA) cycle, with key exit points highlighted in green. Both figures were created with BioRender.com.
    }
    \label{fig:glycolysis_TCA}
\end{figure}

\subsection{Tricarboxylic Acid Cycle}
The TCA cycle, also known as the Krebs cycle or citric acid cycle, represents a central hub of cellular metabolism, integrating multiple catabolic and anabolic pathways. It serves as a key intersection point at which a variety of substrates, including fatty acids, amino acids, and pyruvate, are ultimately funneled into a common metabolic route through their conversion into acetyl-CoA. Acetyl-CoA then is combined with oxaloacetate (\gls{OAA}) to form citrate, marking the entry point of the cycle. 
\\

Citrate is subsequently isomerized to isocitrate, which undergoes two successive oxidative decarboxylation steps. In these reactions, isocitrate is first converted into $\alpha$-ketoglutarate (\gls{aKG}) and then into succinyl-CoA, with the release of two molecules of $CO_2$ and the production of two NADH molecules. Succinyl-CoA is then converted to succinate in a phosphorylation step via succinyl-CoA synthetase  that generates GTP, which can be readily converted into ATP. 
\\

The cycle continues with the oxidation of succinate to fumarate, during which two hydrogen atoms are transferred to FAD, forming $FADH_2$.
This reaction is catalyzed by succinate dehydrogenase, an enzyme that is uniquely embedded in both the TCA cycle and the electron transport chain (\gls{ETC}), linking the two systems directly. Fumarate is then hydrated to malate, which is finally oxidized back to OAA, completing the cycle and allowing it to proceed with another acetyl-CoA molecule. 
\\

Beyond its role in energy production, the TCA cycle functions as a metabolic intersection between catabolic and anabolic processes. Its intermediates are continuously withdrawn for biosynthetic purposes. For example, citrate can be exported to the cytosol, where it is cleaved to generate OAA, a precursor for nucleotide synthesis, and acetyl-CoA, a key building block for lipid synthesis. To maintain cycle continuity despite this metabolic void, several anaplerotic reactions replenish intermediates. Glutamine, for instance, can be converted via glutaminolysis into glutamate and subsequently into
$\alpha$KG, while glucogenic amino acids such as valine and isoleucine can be transaminated and metabolized to yield succinyl-CoA.
\\

In summary, the completion of the TCA cycle generates three NADH, one $FADH_2$, as well as one triphosphate (GTP or ATP), which subsequently feed into the ETC via complex I (CI) and CII to drive OXPHOS and cellular ATP production.
\\

\subsection{Electron Transport Chain and Oxidative Phosphorylation}
The ETC is a series of membrane-embedded protein complexes that transfer electrons through redox reactions, ultimately driving ATP synthesis via OXPHOS. Located in the inner mitochondrial membrane, the ETC couples electron transfer to the translocation of protons across the membrane, establishing the electrochemical gradient (also known as membrane potential) essential for ATP production.

\begin{figure}[H]
    \centering
    \includegraphics[width=\textwidth]{ETC.png}
    \caption{\textbf{Electron transport chain (ETC).}
    The ETC is located in the inner mitochondrial membrane and consists of five complexes(\gls{CI-CV}) that transfer electrons to the final electron acceptor, oxygen. At CIV, two electrons are transferred to reduce ½ $O_2$ to $H_2O$.
    During electron transport, protons are pumped into the intermembrane space, generating a membrane potential that is utilized by Complex V (ATP synthase) to produce ATP. Figure created with BioRender.com.
    }
    \label{fig:ETC}
\end{figure}

Electrons enter the ETC primarily through the reduced cofactors NADH and $FADH_2$,  which are generated by multiple metabolic pathways, including glycolysis, the TCA cycle, and fatty acid oxidation (\gls{FAO}). NADH donates its electrons to Complex I (NADH:ubiquinone oxidoreductase), where they traverse a cascade of iron–sulfur clusters before reducing ubiquinone (coenzyme Q, CoQ) to ubiquinol (\gls{QH2}).
This process is coupled to the translocation of four protons from the mitochondrial matrix into the intermembrane space per NADH molecule. 
\\

In contrast, $FADH_2$ delivers its electrons directly to Complex II (succinate dehydrogenase), which is also a component of the TCA cycle. Electrons from Complex II are likewise transferred to ubiquinone, but unlike Complex I, Complex II does not pump protons. Ubiquinol then shuttles electrons to Complex III (coenzyme Q:cytochrome c oxidoreductase), which further translocates four protons into the intermembrane space per pair of electrons. The electrons are subsequently transferred to the mobile carrier cytochrome c, which transports them to Complex IV (cytochrome c oxidase).
\\

At Complex IV, molecular oxygen acts as the terminal electron acceptor, undergoing reduction to water. This step requires four electrons per $O_2$ molecule and is coupled to the translocation of an additional four protons across the membrane. But since two protons are consumed in water formation, a net of only 2 protons are pumped to the intermembrane space. 
\\

The accumulation of protons in the intermembrane space generates the proton-motive force (\gls{PMF}), which is composed of both a concentration gradient and an electrical membrane potential. This energy is harnessed by CV, through which protons flow back into the mitochondrial matrix. This movement drives the phosphorylation of ADP to ATP, giving rise to the term oxidative phosphorylation. 
\\

In addition to ATP production, the ETC is also a major source of reactive oxygen species (\gls{ROS}).
These arise primarily from the premature leakage of electrons from CI, CII, and CIII, which can partially reduce molecular oxygen, leading to the formation of superoxide and other reactive intermediates.
\\

Importantly, oxidative phosphorylation is substantially more efficient than glycolysis, generating up to \~ 30-32 ATP per glucose molecule, whereas glycolysis alone yields only 2 ATP \citep{mookerjeeQuantifyingIntracellularRates2017}.
However, glycolysis provides a critical advantage under conditions of limited oxygen availability and can be rapidly upregulated to meet acute energy demands. In contrast, although oxidative phosphorylation is comparatively slower, it supports a more sustained and efficient long-term energy supply \citep{melkonianBiochemistryAnaerobicGlycolysis2026}.

\section{Metabolic Dynamics}
Metabolic transitions refer to the dynamic reprogramming of cellular or organismal metabolic pathways in response to changing environmental or physiological conditions \citep{ETCHEGARAYInterplayMetabolismEpigenetics2016}.
These transitions are influenced by both extrinsic and intrinsic factors, including nutrient availability, energy demand, stress responses, and developmental or hormonal cues. 

\subsection{Mechanisms underlying metabolic Transitions}
Rather than functioning as isolated biochemical pathways, metabolism is best understood as a highly interconnected network that integrates cellular requirements with developmental programs. In this context, metabolic activity is continuously adjusted to meet cellular demands, ensuring that energy production and biosynthetic processes are aligned with the cell’s functional state \citep{CHANDELSignalingMetabolism2021,WANGMetaboliteSensingSignaling2018}.
This coordination is essential for supporting processes such as growth, proliferation, and differentiation, while maintaining cellular homeostasis. 

\subsection{Interplay between Metabolism and Gene Expression}
A critical aspect of metabolic regulation lies in its tight interplay with gene expression. While transcriptomic analyses have provided a detailed view of the dynamic gene expression programs underlying processes such as neural development, these programs are ultimately executed through corresponding metabolic changes. Transcriptional states shape metabolic activity by regulating the expression and function of metabolic enzymes. Conversely, metabolites  themselves  and the bioenergetic state of the cell can influence gene expression through epigenetic modifications and signaling mechanisms, establishing a bidirectional relationship between metabolism and transcriptional control (reviewed in \citep{CARTHEWGeneRegulationCellular2021,TARAZONAExploringInfluenceCell2020}).
\\

For instance, intracellular levels of key metabolic intermediates such as acetyl-CoA indirectly promotes histone acetylation, thereby enhancing the transcription of genes involved in cellular growth \citep{CAIAcetylCoAInducesCell2011}.
In addition, metabolic intermediates can function directly as signaling molecules. Lactate, the end product of glycolysis that accumulates under hypoxic conditions, has been shown to promote the stabilization of the protein NDRG3, thereby sustaining the hypoxic response \citep{LEELactateInducedResponseHypoxia2015}.

\subsection{Metabolic Transition from Proliferation to Differentiation}
Large-scale transcriptomic and proteomic analyses, together with genetic perturbation studies, have established metabolic state as a key regulator of neural development \citep{iwataMitochondriaMetabolismSets2023,TARAZONAExploringInfluenceCell2020,ZHENGMetabolicReprogrammingNeuronal2016,OBRIENDifferentiationHumanNeural2015}.
Evidence from both embryonic stem cells and adult neural stem cells consistently points to a close coupling between metabolic phenotype and cell fate decisions.
\\

A central feature of neurogenesis is a progressive metabolic transition along the lineage from stem and progenitor cells to differentiated neurons. This transition reflects changing cellular demands: rapidly proliferating cells require substantial amounts of biosynthetic precursors, including amino acids, lipids, nucleotides, and ATP, to sustain growth and division \citep{ZHANGMetabolicRegulationPluripotent2012,FOLMESSomaticOxidativeBioenergetics2011}.
\\

In general, proliferating cells rely primarily on glycolysis for ATP production, even in the presence of oxygen \citep{METALLOUnderstandingMetabolicRegulation2013}.
In contrast, during differentiation cells transition toward a mitochondrial metabolism dominated by OXPHOS.

\subsection{Limitations of Metabolic Studies}
Although similar metabolic transitions are thought to occur in NPCs, metabolic profiling in these systems remains challenging. Many studies focus on adult neurogenesis and therefore rely on progenitor cells derived from adult brains, which are often expanded \textit{in vitro}.
In addition, adult progenitor cells are predominantly in a quiescent state, a condition that is not present during embryonic neurogenesis. Furthermore, transitions between cell states are continuous rather than discrete, making it difficult to isolate well-defined progenitor populations. Finally, accurately assessing the metabolic state of a cell is challenging because metabolism is highly dynamic. As a result, many studies rely on transcriptomic data, which reflect mRNA levels but do not necessarily correlate with protein abundance or actual metabolic activity.


\subsection{Metabolic State in proliferative neural Stem and Progenitor Cells}
NSCs and NPCs are described as predominantly glycolytic when proliferating. Transcriptomic analyses reveal an enrichment of glycolytic genes, such as hexokinase 2 (\gls{HK2}), GADPH, enolase 1 (\gls{ENO1}), LDHA, and PFK1, which are known to be upregulated by hypoxia-inducible factor-1 (\gls{HIF1}) \citep{LANGEReliefHypoxiaAngiogenesis2016}.
Consistent with this transcriptional profile, proliferating NPCs actively secrete large amounts of lactate, reflecting a strong reliance on aerobic glycolysis \citep{LANGEReliefHypoxiaAngiogenesis2016}.
\\

Low oxygen tension has been shown to support the proliferative state of NSCs \textit{in vitro} by promoting proliferation and self-renewal. Notably, adult NPCs maintain relatively elevated endogenous ROS levels, and perturbation of ROS signaling impairs differentiation \citep{BELLEProliferativeNeuralStem2011},  which at first appears contradictory, as low oxygen availability is generally not associated with increased ROS production. Nevertheless, multiple studies have consistently reported this phenomenon and demonstrated the importance of ROS signaling during neurogenesis \citep{BELLEProliferativeNeuralStem2011,jiUCP2RegulatesEmbryonic2017,chuiOxidativeStressRegulates2020}.
\\

Despite their predominantly  glycolytic phenotype, proliferating NSCs and NPCs remain dependent on functional mitochondria. Genetic disruption of mitochondrial complexes leads to severe developmental consequences. For example, loss of mitochondrial CII results in brain atrophy and early postnatal lethality, primarily due to impaired neuronal maturation and survival, while leaving NSC and progenitor cells largely intact \citep{DIAZ-CASTROResistanceGlialikeCentral2015}.
In contrast, disruption of CI has been shown to strongly impair neurogenesis in neonatal brains, leading to reduced progenitor proliferation and increased apoptosis \citep{CABELLO-RIVERAMitochondrialComplexFunction2019}.
\\

\subsection{Metabolic Reprogramming during Differentiation}
Neural differentiation is accompanied by a marked metabolic shift characterized by increased mitochondrial biogenesis and a progressive downregulation of glycolytic pathways.This metabolic remodeling occurs already at the level of NSCs and NPCs \citep{DONGMetabolicLactateProduction2022} and is maintained throughout neurodevelopment in both embryonic and adult neurogenesis \citep{OBRIENDifferentiationHumanNeural2015,AGOSTINIMetabolicReprogrammingNeuronal2016}.
\\

This transition involves coordinated changes across multiple metabolic pathways. Differentiating cells exhibit increased expression of genes associated with the TCA cycle and FAO, alongside a reduction in lactate secretion, indicating a decreased reliance on glycolytic metabolism and enhanced mitochondrial substrate utilization \citep{LANGEReliefHypoxiaAngiogenesis2016,DONGMetabolicLactateProduction2022,TELLEYSequentialTranscriptionalWaves2016}.
\\

A central feature of this metabolic remodeling is the repression of glycolysis-associated genes, most notably LDHA, a key enzyme supporting aerobic glycolysis.  Reduction of LDHA levels leads to diminished glycolytic flux and a concomitant increase in mitochondrial OXPHOS \citep{FANTINAttenuationLDHAExpression2006}. Sustained LDHA expression during differentiation however, is detrimental to neuronal development and has been linked to neuronal cell death. Moreover, forced maintenance of glycolytic programs, such as overexpression of HK2 or LDHA in NPCs, inhibits neurodifferentiation and biases cells toward astrocytic fate acquisition \citep{ZHENGMetabolicReprogrammingNeuronal2016}. 
\\

Concomitantly, differentiation is associated with increased expression of mitochondrial enzymes, including lactate dehydrogenase B (\gls{LDHB}) and pyruvate dehydrogenase beta \citep{FATHIQuantitativeProteomicsAnalysis2014}.
In contrast to LDHA, LDHB promotes the conversion of lactate to pyruvate, thereby facilitating its entry into mitochondrial oxidative metabolism. This metabolic shift is further supported by pharmacological studies showing that inhibition of mitochondrial respiratory chain CIII  with antimycin A blocks neuronal differentiation and maintains cells in a more progenitor-like state \citep{PEREIRAInhibitionMitochondrialComplex2013}.
\\

Together, these findings indicate that the downregulation of aerobic glycolysis is not merely a consequence of neuronal differentiation but a prerequisite for proper lineage commitment and normal brain development.

\subsection{Metabolism in mature Neurons}
During late neurogenesis and early postnatal development, neurons must migrate radially or tangentially to their final positions in the cortex or subcortical structures, extend their axons over long distances, and branch extensively to finally form functional synapses. These processes impose substantial anabolic demands, requiring the synthesis of lipids necessary for the production of myelin, amino acids for protein and neurotransmitter synthesis, and nucleotides to support membrane expansion, cytoskeletal remodeling, and synaptic vesicle production. Consequently, developing neurons exhibit a metabolic shift from the glycolytic phenotype of NPCs towards oxidative metabolism.
\\

Mature, postmitotic neurons are highly dependent on aerobic mitochondrial metabolism for ATP production \citep{MAGISTRETTICellularPerspectiveBrain2015}, characterized by increased glucose uptake, facilitated by the increase of neuron-specific glucose transporter expression GLUT3 \citep{AGOSTINIMetabolicReprogrammingNeuronal2016}, 
elevated expression of mitochondrial ETC component CV, which show dramatic upregulation upon neuronal differentiation , and a proportional increase in mitochondrial mass accompanying neuronal growth during differentiation \citep{ZHENGMetabolicReprogrammingNeuronal2016}.
In addition, the localization of immobile mitochondria at nascent presynaptic sites has been shown to be required for axon branching \citep{COURCHETTerminalAxonBranching2013},  potentially reflecting a localized metabolic support of branching through the provision of ATP and biosynthetic intermediates necessary for membrane expansion and cytoskeletal remodeling.
\\

Despite the reliance on OXPHOS, mature neurons retain a degree of aerobic glycolysis, which accounts for \~ 10-12\% of glucose utilization in the adult human brain \citep{VAISHNAVIRegionalAerobicGlycolysis2010}. 
Their metabolic state is dynamically regulated in an activity dependent manner, differing between resting and active states (reviewed in \citep{YELLENFuelingThoughtManagement2018}). Furthermore, given the highly polarized and extended morphology of neurons, metabolic processes can be spatially compartmentalized, allowing distinct subcellular regions to adopt different metabolic profiles.

\section{Mitochondria as a central cellular metabolic and bioenergetic Hub}
Mitochondria have long been recognized as the powerhouse of the cell, generating ATP through OXPHOS and the TCA cycle. However, their role extends far beyond energy production, as they also modulate redox signaling, calcium homeostasis, ion homeostasis, apoptosis and epigenetic regulation \citep{chandelMitochondriaSignalingOrganelles2014}.
In the context of embryonal neurogenesis, mitochondria emerge as central hubs that integrate metabolic status with cellular signaling, epigenetic regulation, and fate decisions of NPCs. This multifaceted role makes mitochondria indispensable for understanding the molecular and cellular mechanisms underlying brain development and associated disorders. 

\subsection{Mitochondrial Structure and Physiology}
Mitochondria are double-membraned organelles, consisting of an outer and an inner membrane separated by the intermembrane space. The inner membrane is extensively folded into cristae, which provide the structural platform for the ETC complexes and support efficient ATP production. In addition, mitochondria harbor their own genome, mitochondrial DNA (\gls{mtDNA}),  that encodes 37 genes. These genes include 13 protein-coding genes, which are all essential subunits of the ETC complexes \citep{andersonSequenceOrganizationHuman1981}.
\\

The integrity of mtDNA is therefore critical for maintaining mitochondrial activity, as mutations can impair electron transport, resulting in energy deficits and increased production of ROS. Such mitochondrial dysfunction is associated with a range of diseases, often including neurodevelopmental abnormalities. 

\subsection{Mitochondrial Metabolites and epigenetic Regulation of neural Progenitor Fate}
During embryonal neurogenesis, a metabolic shift occurs from glycolysis to OXPHOS as the neocortex develops and vascularizatNADion improves. Mitochondria orchestrate this metabolic transition, which is essential for sustaining the energy demands of proliferating neural progenitors and differentiating neurons. This metabolic reprogramming directly impacts the availability of mitochondrial metabolites such as acetyl-CoA, $\alpha$KG, and $^+$, that serve as substrates for epigenetic enzymes, thereby reshaping the epigenetic landscape.
\\

NAD$^+$/NADH ratio, set by the balance between glycolytic  and oxidative metabolism, regulates the activity of NAD$^+$-dependent  Enhanced oxidative metabolism increases NAD$^+$ levels, promoting the activity of sirtuin family deacetylases and leading to histone deacetylation, a process associated with reduced stemness and the promotion of differentiation \citep{saundersMiRNAsRegulateSIRT12010,houtkooperSirtuinsRegulatorsMetabolism2012},  whereas a more reduced NADH-state limits sirtuin activity \citep{haigisMammalianSirtuinsBiological2010}.
\\

In contrast, glycolytic metabolism promotes acetyl-CoA availability through the generation of citrate from glucose within the mitochondria, which is exported to the cytosol and subsequently converted into acetyl-CoA by ATP-citrate lyase, located in the cytoplasm and nucleus. Increased acetyl-CoA levels enhances histone acetylation via histone acetyltransferases, supporting an open chromatin state, associated with stem cell maintenance \citep{wellenATPCitrateLyaseLinks2009}.
\\

In parallel, $\alpha$KG, a key intermediate of the TCA cycle, serves as a cofactor for lysine demethylases, driving DNA and histone demethylation \citep{kaelinInfluenceMetabolismEpigenetics2013}.
In comparison to differentiated cells, naive stem cells were found to sustain elevated intracellular  $\alpha$KG levels, which favors demethylation reactions and supports the maintenance of stemness \citep{careyIntracellularAketoglutarateMaintains2015}.

\subsection{Mitochondrial Dynamics}
Mitochondria are dynamic organelles that undergo continuous rounds of fusion and fission, essential for maintaining mitochondrial health, eliminating damaged proteins, and promoting mtDNA complementation. These dynamics regulate metabolic homeostasis and are particularly important during neurogenesis, where mitochondrial morphology changes as cells commit to neuronal lineages \citep{KHACHOMitochondrialDynamicsImpacts2016,IWATAMitochondrialDynamicsPostmitotic2020}.
\\

In uncommitted NSCs, mitochondria typically exhibit an elongated morphology. As cells commit to a progenitor fate, mitochondria become increasingly fragmented, followed by a re-elongation phase as progenitors differentiate into neurons. Proper fusion and fission processes are crucial for maintaining the balance between self-renewal and differentiation as demonstrated by Khacho et al., where loss of fusion through knockout of mitofusin 1 and 2, as well as OPA1, leads to impaired NSC self-renewal, neurogenesis defects and cognitive impairments \citep{KHACHOMitochondrialDynamicsImpacts2016}.

\section{The Scientific Gap: Why Study Mitochondrial Metabolism During Neocortex Development?}
Mitochondria serve as central representatives for cellular metabolism, offering insights into a cell’s health, metabolic state, and functional status. This is particularly true during developmental processes such as neurogenesis, where mitochondrial function is essential to prevent a wide range of neurological disorders. Oxygen availability and the balance between OXPHOS and glycolysis are especially important, as they  shape cellular behavior throughout embryogenesis. Yet, how these metabolic pathways are engaged, and how they adapt to developmental cues remains poorly understood. 
\\

Modeling mitochondrial metabolism provides an initial approach to explore the metabolic states of these transitioning cells. Characterizing these metabolic states provides important insight into the fundamental mechanisms underlying neocortical development. 
